\documentclass[10pt,a4paper]{amsart}
\usepackage{amsmath,amssymb,amsthm,graphicx}
\usepackage{verbatim}
\newenvironment{reference}{\comment}{\endcomment}
%\newenvironment{reference}{}{}
\newenvironment{slogan}{\comment}{\endcomment}
\newenvironment{history}{\comment}{\endcomment}

% For commutative diagrams we use Xy-pic
\usepackage[all]{xy}

% We use 2cell for 2-commutative diagrams.
\xyoption{2cell}
% Theorem environments.
%
\theoremstyle{plain}
\newtheorem{theorem}[subsection]{Theorem}
\newtheorem{proposition}[subsection]{Proposition}
\newtheorem{lemma}[subsection]{Lemma}

\theoremstyle{definition}
\newtheorem{definition}[subsection]{Definition}
\newtheorem{example}[subsection]{Example}
\newtheorem{exercise}[subsection]{Exercise}
\newtheorem{situation}[subsection]{Situation}

\theoremstyle{remark}
\newtheorem{remark}[subsection]{Remark}
\newtheorem{remarks}[subsection]{Remarks}

\numberwithin{equation}{subsection}

% Macros
%
\def\lim{\mathop{\mathrm{lim}}\nolimits}
\def\colim{\mathop{\mathrm{colim}}\nolimits}
\def\Spec{\mathop{\mathrm{Spec}}}
\def\Hom{\mathop{\mathrm{Hom}}\nolimits}
\def\Ext{\mathop{\mathrm{Ext}}\nolimits}
\def\SheafHom{\mathop{\mathcal{H}\!\mathit{om}}\nolimits}
\def\SheafExt{\mathop{\mathcal{E}\!\mathit{xt}}\nolimits}
\def\Sch{\mathit{Sch}}
\def\Mor{\mathop{\mathrm{Mor}}\nolimits}
\def\Ob{\mathop{\mathrm{Ob}}\nolimits}
\def\Sh{\mathop{\mathit{Sh}}\nolimits}
\def\NL{\mathop{N\!L}\nolimits}
\def\CH{\mathop{\mathrm{CH}}\nolimits}
\def\proetale{{pro\text{-}\acute{e}tale}}
\def\etale{{\acute{e}tale}}
\def\QCoh{\mathit{QCoh}}
\def\Ker{\mathop{\mathrm{Ker}}}
\def\Im{\mathop{\mathrm{Im}}}
\def\Coker{\mathop{\mathrm{Coker}}}
\def\Coim{\mathop{\mathrm{Coim}}}

% Boxtimes
%
\DeclareMathSymbol{\boxtimes}{\mathbin}{AMSa}{"02}

%
% Macros for moduli stacks/spaces
%
\def\QCohstack{\mathcal{QC}\!\mathit{oh}}
\def\Cohstack{\mathcal{C}\!\mathit{oh}}
\def\Spacesstack{\mathcal{S}\!\mathit{paces}}
\def\Quotfunctor{\mathrm{Quot}}
\def\Hilbfunctor{\mathrm{Hilb}}
\def\Curvesstack{\mathcal{C}\!\mathit{urves}}
\def\Polarizedstack{\mathcal{P}\!\mathit{olarized}}
\def\Complexesstack{\mathcal{C}\!\mathit{omplexes}}
% \Pic is the operator that assigns to X its picard group, usage \Pic(X)
% \Picardstack_{X/B} denotes the Picard stack of X over B
% \Picardfunctor_{X/B} denotes the Picard functor of X over B
\def\Pic{\mathop{\mathrm{Pic}}\nolimits}
\def\Picardstack{\mathcal{P}\!\mathit{ic}}
\def\Picardfunctor{\mathrm{Pic}}
\def\Deformationcategory{\mathcal{D}\!\mathit{ef}}
\usepackage{cleveref}
\crefname{originalreference}{original reference}{original references}
\Crefname{originalreference}{Original reference}{Original references}
\makeatletter
\newlabel{lemma-H1-injective}{{0AXB}{source}{}{sourcepointer-1}{}}
\newlabel{lemma-H1-injective@cref}{{[originalreference][0][]0AXB}{[1][0][]source}{}{}{}}
\makeatother
\crefname{originalsection}{original section}{original sections}
\Crefname{originalsection}{Original section}{Original sections}
\makeatletter
\newlabel{sites-cohomology-example-hom-complex-into-sheaf}{{07AA}{source}{}{sourcepointer-2}{}}
\newlabel{sites-cohomology-example-hom-complex-into-sheaf@cref}{{[originalsection][0][]07AA}{[1][0][]source}{}{}{}}
\makeatother
\crefname{originalreference}{original reference}{original references}
\Crefname{originalreference}{Original reference}{Original references}
\makeatletter
\newlabel{situation-vanishing}{{07UY}{source}{}{sourcepointer-3}{}}
\newlabel{situation-vanishing@cref}{{[originalreference][0][]07UY}{[1][0][]source}{}{}{}}
\makeatother
\title{Vanishing of maps from the residue field shifted by one into the pushforward of the structure sheaf}
\author{The Stacks Project Authors}
\date{stacks-project:resolve.tex:local-snapshot-20261012; GNU FDL 1.2 or later}
\begin{document}
\maketitle
\noindent\textit{Editorial note.} All delivered mathematical text is the author's own, copied from the pinned Stacks Project chapter file resolve.tex; no statement or proof step was rewritten, reordered or removed.\par\medskip
\noindent\textit{Editorial note.} The counted proof is the author's complete argument: the adjunction translating the map into a map out of the derived pullback, the identification of that map with an extension class through the Ext spectral sequence, and the splitting argument using the lift of the constant section and the normality of the base.\par\medskip
\noindent\textit{Editorial note.} The label lemma-R1-injective names the counted statement, and the proof environment immediately following it in the source is delivered verbatim as its proof.\par\medskip
\noindent\textit{Editorial note.} The counted statement is carried by the author's own lemma environment in the chapter source.\par\medskip
\noindent\textit{Editorial note.} Pre-open gate: python operations/agent-authoring/tools/precard\_gate.py --source resolve.tex --label lemma-R1-injective --start 1756 --end 1804 --local-id 28098 returned RESULT=OK (L28098 free, occupancy FREE, public interval clear). Pre-assembly tag self-check: canonical tag resolved and not present in reference/public-items.jsonl.\par\medskip
\noindent\textit{Editorial note.} The statement and proof are copied from the pinned chapter file; the cited inputs are kept as external pointers with their canonical Stacks tags.\par\medskip
\noindent\textit{Allowed standard prerequisites (editorial).}\begin{itemize}
\item Adjointness of derived pushforward and derived pullback. The standard adjunction identifying maps out of the derived pushforward with maps into the derived pullback of the structure sheaf. Replaces a map from the residue field shifted by one into the derived pushforward by a map from the derived pullback of the residue field shifted by one into the structure sheaf.
\item Cohomology sheaves of the derived pullback of the residue field. The computation of the cohomology sheaves of the derived pullback of the residue field, vanishing above degree one, giving the structure sheaf of the closed fibre in degree one, and in degrees at most zero giving some module over it. Supplies the two ends of the spectral sequence computing the hom space into the structure sheaf.
\item The spectral sequence computing Ext out of a complex with two cohomology sheaves. Cohomology on Sites, the example computing the hom complex into a sheaf. Identifies the hom space in the derived category with the first Ext group out of the structure sheaf of the closed fibre.
\item Splitting of the pullback of the extension along the closed immersion. Same chapter, the lemma asserting that the pullback of a first cohomology class to the open vanishes. Shows that the pullback of the extension class to the preimage of the open is zero, which is what lets the constant section be lifted.
\end{itemize}\medskip
\ifdefined\hypertarget\hypertarget{sourcepointer-1}{}\fi
\noindent\textit{Source reference.} Original reference 0AXB: sources/stacks-resolve/pinned-original-source.tex, source lines 1692--1692. Prerequisite pointer invoked at proof line(s) 1792; it names a standard fact of the same chapter that the counted argument uses, and every step of the counted proof built on it is retained verbatim in the delivered body.\par\medskip
\ifdefined\hypertarget\hypertarget{sourcepointer-2}{}\fi
\noindent\textit{Source reference.} Original section 07AA: sources/stacks-chow/tags.txt, source lines 9300--9300. Prerequisite pointer invoked at proof line(s) 1783; it names a standard fact of a sibling chapter that the counted argument uses, and every step of the counted proof built on it is retained verbatim in the delivered body.\par\medskip
\ifdefined\hypertarget\hypertarget{sourcepointer-3}{}\fi
\noindent\textit{Source reference.} Original reference 07UY: sources/stacks-resolve/pinned-original-source.tex, source lines 1584--1584. Prerequisite pointer invoked at proof line(s) 1758; it names a standard fact of the same chapter that the counted argument uses, and every step of the counted proof built on it is retained verbatim in the delivered body.\par\medskip
\begin{lemma}
\label{lemma-R1-injective}
In Situation \ref{situation-vanishing} assume $X$ is normal and $A$ Nagata.
Then
$$
\Hom_{D(A)}(\kappa[-1], Rf_*\mathcal{O}_X)
$$
is zero. This uses $D(A) = D_\QCoh(\mathcal{O}_S)$ to think of
$Rf_*\mathcal{O}_X$ as an object of $D(A)$.
\end{lemma}

\begin{proof}
By adjointness of $Rf_*$ and $Lf^*$ such a map is the same thing
as a map $\alpha : Lf^*\kappa[-1] \to \mathcal{O}_X$. Note that
$$
H^i(Lf^*\kappa[-1]) =
\left\{
\begin{matrix}
0 & \text{if} & i > 1 \\
\mathcal{O}_{X_s} & \text{if} & i = 1 \\
\text{some }\mathcal{O}_{X_s}\text{-module} & \text{if} & i \leq 0
\end{matrix}
\right.
$$
Since $\Hom(H^0(Lf^*\kappa[-1]), \mathcal{O}_X) = 0$ as $\mathcal{O}_X$
is torsion free, the spectral sequence for $\Ext$
(Cohomology on Sites, Example
\ref{sites-cohomology-example-hom-complex-into-sheaf})
implies that
$\Hom_{D(\mathcal{O}_X)}(Lf^*\kappa[-1], \mathcal{O}_X)$ is equal to
$\Ext^1_{\mathcal{O}_X}(\mathcal{O}_{X_s}, \mathcal{O}_X)$.
We conclude that
$\alpha : Lf^*\kappa[-1] \to \mathcal{O}_X$ is given by an extension
$$
0 \to \mathcal{O}_X \to \mathcal{E} \to \mathcal{O}_{X_s} \to 0
$$
By Lemma \ref{lemma-H1-injective} the pullback of this extension
via the surjection $\mathcal{O}_X \to \mathcal{O}_{X_s}$ is zero
(since this pullback is clearly split over $f^{-1}(U)$).
Thus $1 \in \mathcal{O}_{X_s}$ lifts to a global section $s$ of
$\mathcal{E}$. Multiplying $s$ by the ideal sheaf $\mathcal{I}$
of $X_s$ we obtain an $\mathcal{O}_X$-module map
$c_s : \mathcal{I} \to \mathcal{O}_X$. Applying $f_*$ we obtain
an $A$-linear map $f_*c_s : \mathfrak m \to A$. Since $A$ is
a Noetherian normal local domain this map is given by multiplication
by an element $a \in A$. Changing $s$ into $s -  a$ we find that
$s$ is annihilated by $\mathcal{I}$ and the extension is trivial
as desired.
\end{proof}

\end{document}
